\section{The Hessian spectrum and product exclusion}\label{sec:spectrum}
Clifford operations permute Pauli operators up to phase. The moduli of a state's
Pauli expectations therefore survive every Clifford change of frame and can
obstruct conversion to independent resources. Their fourth moment will give a
product-state obstruction; their rare largest nontrivial values will later
obstruct a short sum-of-cubes synthesis.

For a homogeneous cubic $F$ on $\Fp^k$, put
$|F\rangle=p^{-k/2}\sum_u\omega^{F(u)}|u\rangle$ and
$H_F(v)=(\partial_i\partial_jF)(v)$. We use natural logarithms and define
\[
 M_\alpha(|\psi\rangle)=\frac{1}{1-\alpha}
 \log\left(p^{-k}\sum_{v,w}|\langle\psi|X(v)Z(w)|\psi\rangle|^{2\alpha}\right)
 \quad(\alpha>0,\ \alpha\ne1).
\]
These stabilizer R\'enyi entropies quantify nonstabilizerness, usually called
magic, with the normalization in
\cite{leone,wangli}; they vanish on pure stabilizer states and are additive
under tensor products. The following quadratic-sum calculation makes them
accessible without enumerating all $p^{2k}$ Paulis.

\begin{theorem}[Rank formula]\label{thm:hessian}
\coverage{absent}
Let $p\geq5$, and set $r(v)=\rk H_F(v)$ and
$N_r=\#\{v:r(v)=r\}$. Then
\begin{equation}\label{eq:pauli-spectrum}
 |\langle F|X(v)Z(w)|F\rangle|=
 \begin{cases}p^{-r(v)/2},&w\in\im H_F(v),\\0,&\text{otherwise.}\end{cases}
\end{equation}
Consequently
\begin{equation}\label{eq:rank-entropy}
 M_\alpha(|F\rangle)=\frac{1}{1-\alpha}
 \log\left(p^{-k}\sum_{r=0}^kN_rp^{(1-\alpha)r}\right).
\end{equation}
\end{theorem}
\begin{proof}
The expectation is $p^{-k}\sum_u\omega^{F(u)-F(u+v)+w^Tu}$.
Writing $c_v=\nabla F(v)$, the exponent, up to a constant, is
$-\tfrac12u^TH_F(v)u+(w-c_v)^Tu$.
The sum over the radical vanishes unless $w-c_v\in\im H_F(v)$.
Euler's identity $H_F(v)v=2c_v$ makes this equivalent to the condition
in~\eqref{eq:pauli-spectrum}. On a complementary subspace diagonalize the
quadratic form. Each nonzero quadratic coefficient contributes a Gauss sum
of modulus $\sqrt p$: its squared modulus follows by replacing a pair
$(x,y)$ with $(x-y,x+y)$ in the double sum, an invertible substitution.
The radical contributes $p^{k-r(v)}$, proving the modulus formula.
There are $p^{r(v)}$ admissible $w$ for each $v$; summing their powers proves
\eqref{eq:rank-entropy}.
\proves{thm:hessian}
\end{proof}

Kagamihara--Tsuchiya~\cite[Theorem 1]{kagamihara} give the binary
rank-reduction formula, with $C(x)+C(x)^T$ replacing the odd-characteristic
Hessian. Chen--Yan--Zhou~\cite{chen} supply related hypergraph-state moment
methods. We use the established rank-reduction method in odd characteristic;
the contribution below is the geometric evaluation and its resource consequences.
For example, a single-qudit cubic $F(u)=u^3$ has rank zero at $u=0$ and rank
one at all other points. Thus $N_0=1$, $N_1=p-1$, and
$M_2=\log(p^2/(2p-1))$. Tensor products add these entropies. The same formula
now turns the multivariate rank counts into an exact resource comparison.

\begin{proposition}[Exact conic spectra]\label{prop:spectra}
\coverage{absent}\uses{prop:conic-input,thm:hessian}\evidence{rank-census}
The nonzero entries of the vector rank distributions for the two conic phases are
\[
\begin{array}{c|rrrrr}
p=7:\ r&0&3&4&5&6\\\hline
N_r&1&48&2940&26502&88158
\end{array}
\]
\[
\begin{array}{c|rrrrrrr}
p=11:\ r&0&5&6&7&8&9&10\\\hline
N_r&1&120&14520&80520&21442960&2387485210&23528401270.
\end{array}
\]
In each case the minimum positive rank is $(p-1)/2$, attained on exactly
$p+1$ projective points. In the binary-form coordinates these are $s=0$ and
the perfect $(p-3)$rd powers. In particular,
\[
 M_2(F_7)=\log\frac{1977326743}{78835},\qquad
 M_2(F_{11})=\log\frac{11^{19}}{7151076991}.
\]
\end{proposition}
\begin{proof}[Finite verification and deduction]
The census visits vectors with first nonzero coordinate one. Since
$H_F(av)=aH_F(v)$ for $a\ne0$, every projective representative contributes
$p-1$ vectors; the zero vector is added separately. The domains have
$19608$ and $2593742460$ representatives. The $p=7$ census has an independent
implementation; the $p=11$ exhaustive execution is checked against it in the
smaller case, not independently duplicated at its full size.
Substitution verifies the $p+1$ perfect-power points. The census has exactly
that many minimum-rank lines, so this containment is an equality over $\Fp$.
Equation~\eqref{eq:rank-entropy} gives the entropy values.
\proves{prop:spectra}\uses{thm:hessian}\evidence{rank-census}
\end{proof}

\begin{lemma}[Product ceiling]\label{lem:product-ceiling}
\coverage{absent}
For a pure state in dimension $D$ with an orthogonal unitary Pauli basis,
$M_2\leq\log((D+1)/2)$. Thus a pure product on blocks of $a$ and $k-a$
prime-dimensional qudits satisfies
\begin{equation}\label{eq:product-bound}
 M_2\leq\log\frac{(p^a+1)(p^{k-a}+1)}4.
\end{equation}
\end{lemma}
\begin{proof}
Put $x_P=|\langle P\rangle|^2$. Purity gives $\sum_Px_P=D$ and
$x_I=1$. Cauchy--Schwarz on the $D^2-1$ other entries yields
$\sum_Px_P^2\geq1+(D-1)^2/(D^2-1)=2D/(D+1)$.
Insert this into the definition of $M_2$. For a product the Pauli sum
factorizes, so the entropies add.
\proves{lem:product-ceiling}
\end{proof}
The product-ceiling method has predecessors in the characteristic-function
and qudit entropy literature~\cite{dai,wangli}. The bound need not be attained.

\begin{corollary}[Clifford-product exclusion]\label{cor:product-exclusion}
\coverage{absent}\uses{prop:spectra,lem:product-ceiling}\evidence{magic-comparison}
The ten-qudit state $|F_{11}\rangle$ is not Clifford-equivalent to a product
across any nontrivial bipartition. The same conclusion holds for the dihedral
and Borel six-qudit conic trades over $\mathbb F_7$ in
Appendix~\ref{sec:finite-data}. This entropy test alone does not establish
the conclusion for the six-qudit Clebsch state $|F_7\rangle$.
\end{corollary}
\begin{proof}
For fixed $k$, the largest right side of~\eqref{eq:product-bound} occurs at
$a=1$ or $k-1$: the variable part $p^a+p^{k-a}$ increases toward the
endpoints. At $p=11,k=10$ the bound is $22.6797\ldots$, below
$22.8695\ldots=M_2(F_{11})$. At $p=7,k=6$ it is $10.4228\ldots$;
the other two trades give $10.4246\ldots$ and $10.4840\ldots$, whereas
$M_2(F_7)=10.1299\ldots$. These comparisons are checked as exact rational
inequalities before taking logarithms. Clifford conjugation permutes the
Pauli moduli, so it preserves $M_2$.
\proves{cor:product-exclusion}\uses{prop:spectra,lem:product-ceiling}\evidence{magic-comparison}
\end{proof}

Two disjoint $xyz$ phases at $p=7$ have
$M_2=8.8047\ldots$, also below $F_7$.
The trace cubic $F(u)=\Tr_{\mathbb F_{p^k}/\Fp}(u^3)$ has full Hessian rank
at every nonzero $u$, by nondegeneracy of the trace pairing, and hence
$M_2=\log(D^2/(2D-1))$, $D=p^k$. These finite-field cubic fiducials belong
to the Alltop--MUB construction~\cite{alltop,klappenecker,damski}.
Feng--Luo~\cite[Proposition 2 and Appendix C]{feng} prove single-prime
L1 optimality among diagonal gates. Neither that result nor the L1 ceiling
of~\cite{damski} is an unrestricted $M_2$ maximality theorem.
The conjecture in~\cite{knipfer} concerns two prime-dimensional qudits;
we do not extend its attribution to arbitrary $k$.
